% !PRINTED 35
\weberpage{35}
%<<CONT

\begin{equation}
\begin{aligned}
A&=Q'A'+A'',\\
A'&=Q''A''+A''',\\
&\cdots\\
A^{(\nu-3)}&=Q^{(\nu-2)}A^{(\nu-2)}+A^{(\nu-1)},\\
A^{(\nu-2)}&=Q^{(\nu-1)}A^{(\nu-1)}+A^{(\nu)},
\end{aligned}
\tag{3}
\end{equation}
where $A^{(\nu)}$ is a constant.

If now $A,A'$ have any common divisor, then by the theorems of the preceding paragraph, as is seen from the equations (3), it is also a divisor of $A'',A''',A''''$ and so on down to $A^{(\nu)}$. If $A^{(\nu)}$ is a non-zero constant, then no divisor of $f(x)$ and $\varphi(x)$ depending on $x$ can exist. Such functions are called \gesperrt{prime to one another}, or \gesperrt{relatively prime}. One also says, taking account only of divisors depending on $x$, that the functions have no common divisor.

The condition that the two functions have a common divisor is therefore
\begin{equation}
 A^{(\nu)}=0.
\tag{4}
\end{equation}
In this case $A^{(\nu-1)}$ is a divisor of $A^{(\nu-2)}$, as the last equation (3) shows, and by the preceding equation it is a divisor of $A^{(\nu-3)}$, and so on, hence also a divisor of $A$ and $A'$, that is, of $f$ and $\varphi$. Conversely, every common divisor of $A,A'$ is also a divisor of $A^{(\nu-1)}$; accordingly $A^{(\nu-1)}$ is called the \gesperrt{greatest common divisor} of $f$ and $\varphi$. The algorithm (3) shows that $A^{(\nu)}$ or $A^{(\nu-1)}$ can be derived from the coefficients of $f$ and $\varphi$ by the \gesperrt{rational operations of arithmetic}, and in such a manner that division is always only by the coefficients of the highest powers of $x$ in the functions $A,A',A'',\ldots$ which are not zero.

We shall apply these considerations to two examples.

First let
\begin{equation}
\begin{aligned}
 A&=a_0x^2+a_1x+a_2,\\
 B&=b_0x^2+b_1x+b_2
\end{aligned}
\tag{5}
\end{equation}
be two functions of the second degree, so that $a_0,b_0$ are non-zero.
